\chapter{Results}\label{chap:results}

\section{Dataset Validation Results}\label{sec:results-validation}

This section will report the final validation status of the model-ready
unambiguous-DNA datasets and summarize the effect of cleaning ambiguous symbols.
The completed Bhandari human audit found 162,863 duplicate training rows and
25,703 duplicate test rows in the legacy processed split. A total of 49,060 test
rows, corresponding to 49.87\% of that partition, belonged to exact sequences
also present in training. No identical sequence carried conflicting labels.
Consequently, metrics obtained from the legacy split are retained only as
development and provenance results.

The corrective pipeline produced 33,925 unique strict-DNA promoter windows and
the same number of deterministic mononucleotide shuffle controls. Its 54,280
training rows and 13,570 test rows contain neither internal duplicates nor exact
train--test overlap.

The same audit was then applied to the 26 CDBProm, CNNProm, and GraphPro pairs
that include source-provided non-promoters. Table~\ref{tab:results-real-negative-integrity}
summarizes the outcome by source. No pair passed the strict zero-duplicate and
zero-overlap criterion. Across the three sources, 129,739 of 132,272 test rows
(98.09\%) belong to exact sequences present in training. Four pairs also contain
label conflicts, dominated by 1,345 conflicting sequences in the GraphPro
\textit{E. coli} pair.

\begin{table}[htbp]
  \centering
  \caption{Exact-sequence integrity audit of datasets with source-provided non-promoters. ``Test overlap'' is the percentage of test rows whose normalized sequence occurs in training. None of the audited pairs is eligible for confirmatory evaluation without corrective group-safe partitioning.}
  \label{tab:results-real-negative-integrity}
  \begin{tabular}{lrrrr}
    \toprule
    Source & Pairs & Eligible & Test overlap (\%) & Label conflicts \\
    \midrule
    CDBProm & 11 & 0 & 97.86 & 2 \\
    CNNProm & 7 & 0 & 98.19 & 0 \\
    GraphPro & 8 & 0 & 97.90 & 1346 \\
    \midrule
    Total & 26 & 0 & 98.09 & 1348 \\
    \bottomrule
  \end{tabular}
\end{table}

Provenance tracing showed that the legacy cropping notebook opened its outputs
in append mode. Repeated executions therefore wrote each promoter three times
and each non-promoter four times before row-level partitioning. After changing
the output to overwrite mode and regenerating all 26 pairs, test-row overlap
fell from 98.09\% to 2.17\%, internal duplicate occurrences fell from 379,046
to 1,614, and 9 pairs satisfied the strict integrity criterion without further
correction. The 1,348 unique label conflicts were unchanged, demonstrating that
they are independent of the append defect.

The subsequent mandatory collision filter removed every post-cropping sequence
observed with both labels from both classes. It removed 1,348 positive and 1,348
negative rows, reducing the collection from 183,701 cropped rows to 181,005
rows. Within-class deduplication then removed 301 positive and 759 negative
rows, and the strict-DNA filter removed a further six positive and two negative
rows. The resulting 179,937 unique strict-DNA sequences were assigned
deterministically to train or test as indivisible exact-sequence groups. The
canonical v3 audit contained zero internal duplicates, zero label conflicts,
and zero train--test overlap in all 26 pairs. Its deterministic audit identifier
is \texttt{87371a88ca1c5456db340201de5c19558b1a5773cb441787ffffc0803b66d853}.

Table~\ref{tab:results-filtering-stage-counts} reports the class counts at each
completed construction stage. Partition assignment does not remove rows.

\begin{table}[htbp]
  \centering
  \caption{Aggregate class counts through the corrected CDBProm, CNNProm, and GraphPro construction pipeline. Removed counts are relative to the immediately preceding stage.}
  \label{tab:results-filtering-stage-counts}
  \begin{tabular}{lrrrrr}
    \toprule
    Stage & Positive & Negative & Total & Removed P & Removed N \\
    \midrule
    After cropping & 77,554 & 106,147 & 183,701 & 0 & 0 \\
    After collision removal & 76,206 & 104,799 & 181,005 & 1,348 & 1,348 \\
    After within-class deduplication & 75,905 & 104,040 & 179,945 & 301 & 759 \\
    After strict-DNA filtering & 75,899 & 104,038 & 179,937 & 6 & 2 \\
    Assigned to training & 60,717 & 83,229 & 143,946 & 0 & 0 \\
    Assigned to test & 15,182 & 20,809 & 35,991 & 0 & 0 \\
    \bottomrule
  \end{tabular}
\end{table}

The GraphPro \textit{E. coli} dataset accounts for 1,345 of the 1,348
cross-label collisions. The original positive and negative files already
contain 81 bp sequences, and all 1,345 conflicting sequences have identical
FASTA headers, genomic coordinates, strands, and nucleotide strings in both
classes. Thus, these collisions precede cropping and are not introduced by the
CapsProm pipeline. Table~\ref{tab:results-graphpro-ecoli-filtering} shows their
class-specific effect.

\begin{table}[htbp]
  \centering
  \caption{Class-specific effect of mandatory collision removal on the GraphPro \textit{E. coli} dataset. The raw, pre-crop, and post-crop collision sets are identical.}
  \label{tab:results-graphpro-ecoli-filtering}
  \begin{tabular}{lrrr}
    \toprule
    Class & Raw rows & Collision rows removed & Rows retained \\
    \midrule
    Promoter & 2,519 & 1,345 (53.39\%) & 1,174 \\
    Non-promoter & 3,874 & 1,345 (34.72\%) & 2,529 \\
    \midrule
    Total & 6,393 & 2,690 (42.08\%) & 3,703 \\
    \bottomrule
  \end{tabular}
\end{table}

Exact-identity correction does not exclude highly similar but non-identical
sequences. Table~\ref{tab:results-near-duplicate-sensitivity} therefore reports
the predeclared Hamming-distance sensitivity audit across all 35,991 test
sequences. The broadest threshold found 96 affected test sequences (0.267\%).
All 96 have at least one same-label training neighbor, while none has an
opposite-label neighbor within eight substitutions.

\begin{table}[htbp]
  \centering
  \caption{Train/test near-duplicate sensitivity after canonical v3 exact-sequence correction. Counts report test sequences with at least one training neighbor within each cumulative Hamming-distance threshold.}
  \label{tab:results-near-duplicate-sensitivity}
  \begin{tabular}{rrrr}
    \toprule
    Maximum distance & Minimum identity (\%) & Test sequences & Test share (\%) \\
    \midrule
    1 & 98.77 & 17 & 0.047 \\
    2 & 97.53 & 32 & 0.089 \\
    4 & 95.06 & 49 & 0.136 \\
    8 & 90.12 & 96 & 0.267 \\
    \bottomrule
  \end{tabular}
\end{table}

\section{Baseline Results}\label{sec:results-baselines}

The first leakage-controlled runs compare logistic regression and XGBoost using
the same 66-dimensional representation: normalized 3-mer frequencies, global
GC content, and GC skew. Table~\ref{tab:results-bhandari-shuffle-baselines}
reports the fixed-seed development results. These values characterize the
Bhandari shuffle-control task and must not be interpreted as performance against
genomic non-promoters.

\begin{table}[htbp]
  \centering
  \caption{Fixed-seed development results on the deduplicated Bhandari human shuffle-control dataset. XGBoost improves all reported effectiveness and calibration measures relative to logistic regression, but the synthetic negative class limits the conclusion to this controlled task.}
  \label{tab:results-bhandari-shuffle-baselines}
  \begin{tabular}{lrr}
    \toprule
    Metric & Logistic regression & XGBoost \\
    \midrule
    Accuracy & 0.7913 & 0.8211 \\
    Precision & 0.7924 & 0.8291 \\
    Recall & 0.7894 & 0.8088 \\
    Specificity & 0.7932 & 0.8333 \\
    F1 & 0.7909 & 0.8189 \\
    ROC-AUC & 0.8727 & 0.9049 \\
    Average precision & 0.8744 & 0.9102 \\
    MCC & 0.5826 & 0.6423 \\
    Brier score & 0.1486 & 0.1258 \\
    Log-loss & 0.4608 & 0.3951 \\
    Expected calibration error & 0.0572 & 0.0311 \\
    \bottomrule
  \end{tabular}
\end{table}

The observed MCC difference is 0.0597 in favor of XGBoost. Because these are
single-seed development runs and no paired interval or randomization test has
yet been executed, the difference is descriptive rather than confirmatory.

\subsection{Exploratory biological-negative smoke test}

The canonical v3 integration was verified on the GraphPro
\textit{C. phytofermentans} pair using the same 66-dimensional representation.
This test exercised YAML loading, fold-manifest generation, MLflow tracking,
row-level predictions, effectiveness metrics, and calibration metrics. Because
the test outcomes have been observed, Table~\ref{tab:results-real-negative-smoke}
is an exploratory pipeline check and is excluded from confirmatory inference.

\begin{table}[htbp]
  \centering
  \caption{Exploratory fixed-seed smoke-test results on the canonical GraphPro \textit{C. phytofermentans} pair. These values validate pipeline integration and must not be used for confirmatory model selection.}
  \label{tab:results-real-negative-smoke}
  \begin{tabular}{lrr}
    \toprule
    Metric & Logistic regression & XGBoost \\
    \midrule
    MCC & 0.1158 & -0.0245 \\
    ROC-AUC & 0.5748 & 0.5091 \\
    Average precision & 0.5940 & 0.5759 \\
    Balanced accuracy & 0.5578 & 0.4880 \\
    Brier score & 0.2466 & 0.2866 \\
    Log-loss & 0.6864 & 0.7846 \\
    \bottomrule
  \end{tabular}
\end{table}

\section{Feature Engineering Results}\label{sec:results-features}

\pendingresults

\section{CapsProm Results}\label{sec:results-capsprom}

\pendingresults

\section{A-CapsProm Results}\label{sec:results-acapsprom}

\pendingresults
